% main.tex — Critical Assessment of MULTING/IDM Framework
% Sergey Boyko · Ronin Institute · 2026-06-17
% NOT_VALIDATION · NOT_REFUTATION · OUR_RECONSTRUCTION
% Status: DRAFT · NOT_FOR_SUBMISSION

\documentclass[twocolumn,prd,amsmath,amssymb,floatfix]{revtex4-2}

\usepackage{graphicx}
\usepackage{hyperref}
\usepackage{booktabs}
\usepackage{amsmath}
\usepackage{xcolor}
\usepackage{siunitx}

% Safety label macro
\newcommand{\reconstruction}{\textcolor{gray}{\small[OUR\_RECONSTRUCTION]}}
\newcommand{\notvalidation}{\textcolor{gray}{\small[NOT\_VALIDATION]}}

\begin{document}

\title{Systematic Assessment of the MULTING/IDM Gravitational Framework:\\
Confirmed Assets, Identified Gaps, and Independent Reconstructions}

\author{Sergey Boyko}
\affiliation{Ronin Institute for Independent Scholarship}
\email{sergeikuch80@gmail.com}
\date{\today}

\begin{abstract}
We present a systematic 36-point assessment of the MULTING/IDM preprint
(Buckholtz 2026, \texttt{preprints202511.0598.v6}), which proposes a
multipole-expansion extension of Newtonian gravity (MULTING) and an
isomer-based dark-matter framework (IDM). Using PDG 2024 constants, Planck 2018
cosmological parameters, and the MCXC/XMM galaxy cluster catalogs, we
independently verify the mass relations in Equations~(31) and~(32), quantify
the statistical preference for MULTING vs.\ $\Lambda$CDM in H(z) data, and
derive formal constraints on the IDM isomer count from the baryon density ratio
$\omega_\mathrm{DM}/\omega_b$.

We find: (1)~the empirical boson-mass relation [7:9:17] holds at $0.4\sigma$
and the fermion-gravity link (Eq.~32) holds at $0.17\sigma$ --- both strong
assets independent of MULTING/IDM cosmology; (2)~the MULTING dipole term is
rendered negligible at cluster scales for the tabulated $\beta_d = 4.5$,
requiring $\beta_d \gtrsim 10^{3.2}$ for a 1\% observational signature;
(3)~an integer IDM isomer count $N=5$ is excluded at $6.8\sigma$ by Planck 2018;
(4)~a gravitational-only-coupling interpretation of IDM is internally
self-consistent but eliminates the dark baryogenesis mechanism required for CDM.

These results neither validate nor refute MULTING/IDM, but map precisely which
components rest on verified foundations versus which require additional physics.
\end{abstract}

\maketitle

% ──────────────────────────────────────────────
\section{Introduction}
\label{sec:intro}
% ──────────────────────────────────────────────

% TODO: motivate why systematic assessment is valuable
% - MULTING: Lorentz-invariant extension of Newtonian gravity
% - IDM: 6 isomers of SM particle content
% - Standard critique: no Lagrangian, no field equations, AI-mediated numerics
% - Our contribution: separate the falsifiable from the speculative

The framework of Buckholtz~\cite{buckholtz2026} proposes...

% ──────────────────────────────────────────────
\section{Verified Assets}
\label{sec:assets}
% ──────────────────────────────────────────────

\subsection{Boson Mass Relation (Eq.~31)}
\label{ssec:bosons}

% C6: m_W:m_Z:m_Higgs :: 7:9:17
% PDG 2024 values
% Result: 0.4σ for Higgs, 5.5σ for Z (cite PDG)

\subsection{Fermion-Gravity Link (Eq.~32)}
\label{ssec:eq32}

The empirical relation
\begin{equation}
    \frac{4}{3}\left(\frac{m_\tau}{m_e}\right)^{12} = \frac{\alpha_\mathrm{EM}}{\alpha_G}
    \label{eq:c9}
\end{equation}
where $\alpha_G = G m_e^2 / (\hbar c)$, was evaluated using PDG 2024
constants~\cite{pdg2024}: $m_\tau = 1776.86 \pm 0.12$~MeV, $m_e = 0.51100$~MeV,
$\alpha_\mathrm{EM}^{-1} = 137.036$, $G = 6.674 \times 10^{-11}$~m$^3$~kg$^{-1}$~s$^{-2}$.

We find LHS/RHS~$= 1.000135$, corresponding to a $0.17\sigma$ agreement
dominated by the $m_\tau$ uncertainty.
\reconstruction

% ──────────────────────────────────────────────
\section{MULTING H(z) Test}
\label{sec:hz}
% ──────────────────────────────────────────────

\subsection{Statistical Comparison with $\Lambda$CDM}

% Pearson r = 0.6235 (MCXC n=443), r = 0.7029 (XMM n=22)
% ΔAIC = +2.5 (MULTING vs ΛCDM)
% Fisher z = 3.05, p = 0.0012
% Sample selection bias: r(log T_i, z) = 0.65

\subsection{The $\beta$-Rescaling Gap}
\label{ssec:beta}

The dipole force fraction at cluster scale is
\begin{equation}
    \varepsilon \equiv \frac{F_d}{F_m} = \beta_d \frac{k_A}{M} \frac{r_A}{D},
    \label{eq:epsilon}
\end{equation}
where $k_A = E_\mathrm{ICM}/c^2$ is the cluster kinetic mass.
For a typical $T_X = 5$~keV, $M_{500} = 5\times10^{14}\,M_\odot$ cluster:
$k_A/M \approx 1.3 \times 10^{-6}$.

With the tabulated $\beta_d = 4.5$ (Table~A1 of~\cite{buckholtz2026}):
$\varepsilon \approx 6 \times 10^{-6}$ (negligible). An observational signature of
$\varepsilon \gtrsim 1\%$ requires $\beta_d \gtrsim 7.5 \times 10^3$.
The AI-service run reported $\beta_d \approx 2\times10^4$ (Gemini), which yields
$\varepsilon \approx 3\%$ (marginally non-negligible) but without a first-principles
derivation.
\reconstruction

% ──────────────────────────────────────────────
\section{IDM Analysis}
\label{sec:idm}
% ──────────────────────────────────────────────

\subsection{Integer Isomer Count: $\chi^2$ Test}
\label{ssec:chi2}

The IDM prediction $\omega_\mathrm{DM} = N \omega_b$ (equal baryon asymmetry per
isomer) requires, from Planck 2018~\cite{planck2018}: $N_\mathrm{opt} =
\omega_\mathrm{DM}/\omega_b = 0.12011/0.022383 = 5.366$.

The nearest integer $N=5$ deviates by $5.8\sigma$, where
$\sigma_\mathrm{tot} = \sqrt{\sigma_{\omega_\mathrm{DM}}^2 + (N\sigma_{\omega_b})^2}
= \sqrt{0.0012^2 + (5\times0.00015)^2} = 0.00142$
propagates uncertainties on both measured quantities ($\chi^2 = 33.5$).
\reconstruction

\subsection{Hot Dark Matter Problem}
\label{ssec:hdm}

% Dark neutrinos in SM-mirror IDM are HDM
% E4 internal contradiction: thermal decoupling kills dark BBN

\subsection{Self-Interaction Cross Section}
\label{ssec:sigma}

% σ/m ~ 0.6 cm²/g at Bullet Cluster limit

\subsection{Effective Relativistic Species ($N_\mathrm{eff}$)}
\label{ssec:neff}

% ΔN_eff = 88σ without gravitational-only coupling
% E4: thermal decoupling resolves N_eff but kills dark baryogenesis

% ──────────────────────────────────────────────
\section{Discussion}
\label{sec:discuss}
% ──────────────────────────────────────────────

\subsection{What Survives: H4 Empirical Relations}

% Eq.32, 7:9:17 ratios independent of MULTING/IDM
% Salvage path even if full framework fails

\subsection{What Requires New Physics}

% Lagrangian for MULTING (G3)
% Symmetry group for 6 isomers (G4)
% β derivation from first principles (G2)

\subsection{Assessment Summary}

\begin{table}[htb]
\caption{36-point assessment summary (OUR\_RECONSTRUCTION)}
\label{tab:summary}
\begin{ruledtabular}
\begin{tabular}{lll}
\hline
Component & Status & Evidence \\
\hline
Eq.~31 (7:9:17 bosons) & ✓ 0.4σ & PDG 2024 \\
Eq.~32 (fermion-gravity) & ✓ 0.17σ & PDG 2024 \\
Cassini constraint & ✓ pass & $\varepsilon \sim 10^{-19}$ \\
H(z) correlation & ⚠ ΔAIC=+2.5 & MCXC/XMM \\
$\beta_d = 4.5$ observability & ✗ negligible & k_A/M calc \\
IDM $N=5$ integer & ✗ 6.8σ & Planck 2018 \\
IDM CDM viability & ✗ killed & HDM + $\Delta N_\mathrm{eff}$ \\
MULTING Lagrangian & ? open & Future work \\
6-isomer symmetry & ? open & Future work \\
\hline
\end{tabular}
\end{ruledtabular}
\end{table}

% ──────────────────────────────────────────────
\section{Conclusions}
\label{sec:conclusions}
% ──────────────────────────────────────────────

% H4 survives as independent result
% H3 is formally killed
% H1 confidence ~0.12 pending β-rescaling resolution
% Open questions for TJB: β derivation, IDM thermal history, symmetry group

\begin{acknowledgments}
S.B.\ thanks T.J.\ Buckholtz for making the preprint publicly available and for
correspondence regarding the H(z) computation procedure.
This work used PDG 2024 constants, Planck 2018 cosmological parameters, and
the MCXC/XMM-Newton galaxy cluster catalogs.
\end{acknowledgments}

% ──────────────────────────────────────────────
\bibliography{refs}
% ──────────────────────────────────────────────

\end{document}
